\section{Relative radial persistence of the original clearance source}
\label{sec:v68-radial-persistence}
We now restore the smooth physical guard. Analyticity is used only for
the finite sign geometry, not asserted for the guard. In particular an
infinitely flat guard at a zero center value is not assigned a finite
birth order.

Let $D_z^\varepsilon$ be the original first-clearance multiplier of
Section~\ref{sec:v66-physical-angular}. It satisfies
$0\le D_z^\varepsilon\le1$ and the Lipschitz estimate
\eqref{eq:v66-guard-lipschitz}. Write
\begin{equation}\label{eq:v68-positive-center-guard}
 g_z=D_z^\varepsilon(0),\qquad
 L_g=C_g\varepsilon^{-1},\qquad
 k_z^{\rm g}=L_g/g_z\quad\hbox{when }g_z>0.
\end{equation}
At a competing-clearance seam $g_z=1$. A non-seam critical center
with $g_z>0$ is allowed as well. Define the exact weighted fraction
\begin{equation}\label{eq:v68-weighted-radial-fraction}
 a_{z,\lambda}^{\varepsilon}(r)=\frac1{2\pi}\int_0^{2\pi}
 B_{z,\lambda}(r\omega_\phi)D_z^\varepsilon(r\omega_\phi)\,d\phi.
\end{equation}
Thus this is $A_{z,n,k}^{\varepsilon}(r^2/2)/(2\pi)$, not a new
source or a substituted mask.

For a nonzero angular germ and $g_z>0$ set
\begin{align}
 c_{z,\lambda}&=g_z\eta_{z,\lambda},\qquad
 \gamma_{z,\lambda}=J_zc_{z,\lambda},\notag\\
 \mathfrak p_{z,\lambda}
    &=\max\bigl\{(r_{z,\lambda}^*)^{-1},
             2(k_{z,\lambda}^{\rm a}+k_z^{\rm g})\bigr\}.
                                                        \label{eq:v68-persistence-number}
\end{align}
All these numbers are finite at the fixed word. A change of positive
scale of a primitive scalar margin changes none of its angular sets.
The profile certificate is recorded on those exact sets. No finite
value of $\mathfrak p$ is assigned when $g_z=0$ or the angular germ
is identically zero.

\begin{theorem}[Relative radial profile at every finite birth order]
\label{thm:v68-relative-radial-profile}
For $\mathfrak p_{z,\lambda}\le K$ and $0<r<K^{-1}$,
\begin{equation}\label{eq:v68-relative-radial-profile}
 (1-Kr)c_{z,\lambda}r^{d_{z,\lambda}}
 \le a_{z,\lambda}^{\varepsilon}(r)
 \le(1+Kr)c_{z,\lambda}r^{d_{z,\lambda}}.
\end{equation}
Consequently, with $u=r^2/2$, on the common coarea representative,
\begin{equation}\label{eq:v68-relative-density-profile}
 e^{-K_\chi r}(1-Kr)\gamma_{z,\lambda}r^{d_{z,\lambda}}
 \le b_{z,\lambda}^{\varepsilon,\mathrm{clr}}(t_z+u)
 \le e^{K_\chi r}(1+Kr)\gamma_{z,\lambda}r^{d_{z,\lambda}}.
\end{equation}
The constants do not depend on $d_{z,\lambda}$. For a seam and
$d_{z,\lambda}=0$, $\gamma_{z,\lambda}=\beta_{z,\lambda}$.
For $d_{z,\lambda}>0$, $\gamma$ is a leading radial coefficient,
not a nonzero zero-width atom.
\end{theorem}
\begin{proof}
On the circle of radius $r$, the Lipschitz estimate gives
$g_z(1-k_z^{\rm g}r)\le D_z^\varepsilon\le
 g_z(1+k_z^{\rm g}r)$. Multiply by the nonnegative exact indicator
and integrate. Combine this with
\eqref{eq:v68-relative-angular-type}. Put
$x=k_{z,\lambda}^{\rm a}r$ and $y=k_z^{\rm g}r$.
Since $x+y\le Kr/2<1/2$, the lower product is at least
$1-x-y\ge1-Kr$. For the upper product,
$x+y+xy\le2(x+y)\le Kr$. This proves
\eqref{eq:v68-relative-radial-profile}. Multiplication by the unchanged
positive density distortion in
\eqref{eq:v66-multiplicative-profile} proves
\eqref{eq:v68-relative-density-profile}. There is no summation over
nonphysical labels or words in this argument. The last two statements
follow from the values of $g_z$, $d$ and $\eta$.
\end{proof}

\subsection{An outer annulus controls every inner level}
For a fixed $H>0$ with $2H<h_\chi$ and $2K\sqrt H<1$ define
\begin{align}
 q_{z,2H}&=\frac1{2H}\int_0^{2H}
              b_{z,\lambda}^{\varepsilon,\mathrm{clr}}(t_z+v)\,dv,
                \notag\\
 \overline q_{z,H}&=\frac1H\int_H^{2H}
              b_{z,\lambda}^{\varepsilon,\mathrm{clr}}(t_z+v)\,dv,
                                                        \label{eq:v68-annular-averages}\\
 F_{\chi,K}(H)&=
 e^{(2+\sqrt2)K_\chi\sqrt H}
             \frac{1+\sqrt2K\sqrt H}{1-2K\sqrt H}.
                                                        \label{eq:v68-annular-factor}
\end{align}
All coefficients use the original source at its actual roof value.

\begin{proposition}[Order-free positive annular comparison]
\label{prop:v68-annular-comparison}
For every selected pair with $\mathfrak p_{z,\lambda}\le K$,
\begin{align}
 \gamma_{z,\lambda}(2H)^{d_{z,\lambda}/2}
   &\le\frac{e^{2K_\chi\sqrt H}}{1-2K\sqrt H}
                       \overline q_{z,H}
    \le\frac{2e^{2K_\chi\sqrt H}}{1-2K\sqrt H}q_{z,2H},
                                                        \label{eq:v68-mode-annulus}\\
 \operatorname*{ess\,sup}_{0<u<H}
             b_{z,\lambda}^{\varepsilon,\mathrm{clr}}(t_z+u)
   &\le F_{\chi,K}(H)\overline q_{z,H}
    \le2F_{\chi,K}(H)q_{z,2H}.             \label{eq:v68-inner-annulus}
\end{align}
In particular there is no cost growing with the finite birth order.
\end{proposition}
\begin{proof}
For $H<v<2H$, the radial variable lies between $\sqrt{2H}$ and
$2\sqrt H$. Since $d\ge0$, the lower bound in
\eqref{eq:v68-relative-density-profile} is at least
\[
 e^{-2K_\chi\sqrt H}(1-2K\sqrt H)
                              \gamma_{z,\lambda}(2H)^{d/2}.
\]
Average over an interval of length $H$. Positivity gives
$\overline q_{z,H}\le2q_{z,2H}$ and proves
\eqref{eq:v68-mode-annulus}. For $0<u<H$ the upper profile is at
most
$e^{\sqrt2K_\chi\sqrt H}(1+\sqrt2K\sqrt H)
 \gamma_{z,\lambda}(2H)^{d/2}$.
Combine the two inequalities. Only the monotonicity of $r^d$ for
$d\ge0$ was used. In particular neither a factor $2^d$ nor a
maximum allowed value of $d$ is introduced.
\end{proof}

\paragraph{The two excluded zero cases.}
If the angular germ is identically zero, the angular source vanishes
on its certified smaller disk, but source in the rest of the chart
is retained. If $g_z=0$, the guard can be nonzero off the center and
need not be analytic. Neither a division by $g_z$ nor a finite-order
classification of that guarded source is allowed. Such source will
appear in the positive complement. These exclusions do not remove a
zero-tangent-fraction label whose angular germ is nonzero and whose
center guard is positive; all its finite orders are included above.
